% Complete finite-scope proofs for P15-U1.
\section{Finite probability: proofs}
\label{sec:annex-p15-finite-probability}

All outcome spaces in these proofs are finite, and every random variable
has finite range. Sums over value sets use those ranges or finite supersets.
The definitions are those of Chapter~\ref{sec:p15-finite-probability}. The arguments supply the
finite versions of the source's probability, conditioning, independence and
expectation results, with all conditioning denominators stated explicitly
\cite[secs.~1.2, 4.1 and 6.1]{GrinsteadSnellChance2006}.

\begin{annexitem}{annex:p15.finite-laws}{thm:p15-finite-probability-laws}{Finite probability identities and atomic representation}
\begin{proof}
An empty sum gives \(\mathbb P(\varnothing)=0\), and normalization of
\(p\) gives \(\mathbb P(\Omega)=1\). Nonnegativity of each summand
gives \(\mathbb P(A)\geq0\). If \(A\subseteq B\), the sum for
\(B\) contains the sum for \(A\) and only additional nonnegative
terms, proving monotonicity. Taking \(B=\Omega\) gives the upper bound.

If \(A_1,\ldots,A_m\) are pairwise disjoint, each outcome in their
union contributes its mass once to each side of
\[
 \sum_{\omega\in\bigcup_i A_i}p(\omega)
 =\sum_{i=1}^m\sum_{\omega\in A_i}p(\omega).
\]
For \(m=0\), the union and the sum are both empty. Additivity applied
to \(\Omega=A\cup A^c\) proves the complement formula.
The disjoint decompositions
\[
\begin{aligned}
 A\cup B&=(A\setminus B)\cup(B\setminus A)\cup(A\cap B),\\
 A&=(A\setminus B)\cup(A\cap B),\\
 B&=(B\setminus A)\cup(A\cap B)
\end{aligned}
\]
give the two-event union formula after subtracting the duplicated
intersection term.

A finite sum of nonnegative real numbers equals zero exactly when every
summand equals zero. Applying this to \(A\) proves the null-event
criterion. The complement formula then says that \(A\) has probability
one exactly when every point outside it has mass zero, which is equivalent
to containing the support.

For the converse, write \(Q\) for the given nonnegative, normalized,
disjoint-pair-additive function on \(2^\Omega\). Since the empty set is
disjoint from itself, \(Q(\varnothing)=2Q(\varnothing)\), so
\(Q(\varnothing)=0\). Induction on the number of disjoint sets extends
pair additivity to finite additivity. Define \(p(\omega)=Q(\{\omega\})\).
These values are nonnegative, and the singleton partition of \(\Omega\)
shows that they sum to one, hence each is at most one. The singleton
partition of any event \(A\) now gives
\(Q(A)=\sum_{\omega\in A}p(\omega)\). Any mass function inducing
\(Q\) must assign \(Q(\{\omega\})\) to \(\omega\), proving uniqueness.
\end{proof}
\end{annexitem}

\begin{annexitem}{annex:p15.finite-law-projection}{thm:p15-finite-law-projection}{Grouping outcomes and marginalization}
\begin{proof}
The fibers \(\{X=x\}\), as \(x\) ranges over \(S\), are pairwise
disjoint and cover \(\Omega\). Some may be empty; their probability is
then zero. Finite additivity shows that their nonnegative probabilities
sum to one. The union of the fibers with \(x\in D\) is exactly
\(\{X\in D\}\), proving the first formula.

Apply the same argument to the function \((X,Y)\) to obtain a joint mass
function. For a fixed \(x\), the events
\(\{X=x\}\cap\{Y=y\}\), for \(y\in T\), partition
\(\{X=x\}\) up to empty fibers. Their probabilities therefore sum to
\(p_X(x)\). Fixing \(y\) instead proves the other marginal identity.
Finally,
\[
 \{f(X)=u\}=\bigcup_{x\in S:\,f(x)=u}\{X=x\}
\]
is a disjoint union. Summing its masses proves the grouping formula,
including the empty-fiber case.
\end{proof}
\end{annexitem}

\begin{annexitem}{annex:p15.finite-conditioning}{thm:p15-finite-conditioning-laws}{Conditional reweighting and chain rules}
\begin{proof}
For \(\mathbb P(B)>0\), all values of \(p_B\) are nonnegative and
\[
 \sum_{\omega\in\Omega}p_B(\omega)
 =\frac{\sum_{\omega\in B}p(\omega)}{\mathbb P(B)}=1.
\]
Its event probability at \(A\) is
\[
 \sum_{\omega\in A}p_B(\omega)
 =\frac{\sum_{\omega\in A\cap B}p(\omega)}{\mathbb P(B)}
 =\mathbb P(A\mid B).
\]
Thus the finite probability identities apply to \(\mathbb P_B\), and
substitution of \(A=B\) gives probability one. Multiplication of the
defining ratio by its positive denominator proves the multiplication rule.

For each \(1\leq i\leq n\), the prefix event \(B_{i-1}\) has
positive probability by hypothesis, so
\[
 \mathbb P(B_i)=\mathbb P(A_i\mid B_{i-1})\mathbb P(B_{i-1}).
\]
Substitute these equations successively, starting from
\(\mathbb P(B_0)=1\). This gives the chain rule. The final intersection
may have probability zero; it is not used as a conditioning denominator.
If an earlier prefix is null, this displayed conditional product is not
asserted.

If \(\mathbb P(B\cap C)>0\), monotonicity gives
\(\mathbb P(B)>0\), and
\(\mathbb P_B(C)=\mathbb P(B\cap C)/\mathbb P(B)>0\). Hence
both successive conditionings are defined, and
\[
 \mathbb P_B(A\mid C)
 =\frac{\mathbb P_B(A\cap C)}{\mathbb P_B(C)}
 =\frac{\mathbb P(A\cap B\cap C)}{\mathbb P(B\cap C)}
 =\mathbb P(A\mid B\cap C).
\]
\end{proof}
\end{annexitem}

\begin{annexitem}{annex:p15.finite-bayes}{thm:p15-finite-bayes}{Total probability, posterior weights and odds}
\begin{proof}
The events \(E\cap H_i\) are pairwise disjoint and cover \(E\).
Finite additivity therefore gives the first sum in the total-probability
identity. If \(i\notin J\), then \(\mathbb P(H_i)=0\), so
monotonicity gives \(\mathbb P(E\cap H_i)=0\). If \(i\in J\),
the multiplication rule gives
\(\mathbb P(E\cap H_i)=\pi_i\ell_i\). Separating these two cases
proves the second sum without ever forming a conditional probability on
a null hypothesis.

When \(\mathbb P(E)>0\), divide each intersection probability by this
number. This proves the two branches of Bayes' formula. Summing its
numerators gives \(\mathbb P(E)\), so the posterior weights sum to
one. The zero branch also proves preservation of a null hypothesis.

In the two-hypothesis case, the assumptions ensure that both likelihoods
are defined and \((1-\pi)q_0>0\). Thus
\(\mathbb P(E)=\pi q_1+(1-\pi)q_0>0\), and the posterior weight of
\(H^c\) is positive. Dividing the two Bayes formulas cancels their
common denominator and yields
\[
 \frac{\pi q_1}{(1-\pi)q_0}
 =\frac{\pi}{1-\pi}\frac{q_1}{q_0}.
\]
This also covers \(q_1=0\). The statement imposes \(q_0>0\) rather
than assigning an ordinary real value to an infinite likelihood ratio.
\end{proof}
\end{annexitem}

\begin{annexitem}{annex:p15.finite-independence}{thm:p15-finite-independence-laws}{Independence criteria and product models}
\begin{proof}
When \(\mathbb P(B)>0\), dividing the product identity by
\(\mathbb P(B)\) gives \(\mathbb P(A\mid B)=\mathbb P(A)\).
Multiplying this conditional identity by the same denominator proves its
converse. If \(\mathbb P(A)=0\), then \(A\cap B\subseteq A\)
forces \(\mathbb P(A\cap B)=0=\mathbb P(A)\mathbb P(B)\).
Symmetry gives the claim for either null event.

If \(A,B\) are independent, decompose \(B\) into its disjoint parts
\(A\cap B\) and \(A^c\cap B\). Then
\[
 \mathbb P(A^c\cap B)
 =\mathbb P(B)-\mathbb P(A\cap B)
 =(1-\mathbb P(A))\mathbb P(B)
 =\mathbb P(A^c)\mathbb P(B).
\]
Interchanging the events proves independence of \(A,B^c\), and applying
the result once more gives independence of \(A^c,B^c\). No division is
used here, so zero and one probabilities are included.

Suppose the joint law factors. For any \(D\subseteq S,F\subseteq T\),
the relevant joint fibers are disjoint, so
\[
\begin{aligned}
 \mathbb P(\{X\in D\}\cap\{Y\in F\})
 &=\sum_{x\in D}\sum_{y\in F}p_{X,Y}(x,y)\\
 &=\sum_{x\in D}\sum_{y\in F}p_X(x)p_Y(y)\\
 &=\left(\sum_{x\in D}p_X(x)\right)
   \left(\sum_{y\in F}p_Y(y)\right).
\end{aligned}
\]
The last line is \(\mathbb P(X\in D)\mathbb P(Y\in F)\).
Conversely, apply event independence to the singleton sets
\(D=\{x\}\), \(F=\{y\}\) to obtain the factorization of every
joint mass.

For the product construction, nonnegativity is immediate and
\[
 \sum_{(x,y)\in S\times T}a(x)b(y)
 =\left(\sum_{x\in S}a(x)\right)
  \left(\sum_{y\in T}b(y)\right)=1.
\]
Summing over one coordinate gives respectively \(a(x)\) and \(b(y)\)
as the marginal masses. Their product is the assigned joint mass, proving
independence of the coordinate variables. This establishes a model with
independent coordinates; it does not infer independence of observed
quantities from knowledge of their marginal laws alone.
\end{proof}
\end{annexitem}

\begin{annexitem}{annex:p15.finite-expectation}{thm:p15-finite-expectation-laws}{Finite expectation identities}
\begin{proof}
Group the finite outcome sum by the fibers of \(X\):
\[
\begin{aligned}
 \mathbb E[f(X)]
 &=\sum_{x\in S}\sum_{\omega:X(\omega)=x}f(x)p(\omega)\\
 &=\sum_{x\in S}f(x)\mathbb P(X=x).
\end{aligned}
\]
Empty fibers contribute zero. For \(\mathbf1_A\), the defining sum
contains \(p(\omega)\) exactly at the outcomes of \(A\), hence equals
\(\mathbb P(A)\).

Pointwise arithmetic and finite distributivity give
\[
\begin{aligned}
 \mathbb E\left[c_0+\sum_{i=1}^n c_iX_i\right]
 &=\sum_{\omega\in\Omega}
      \left(c_0+\sum_{i=1}^n c_iX_i(\omega)\right)p(\omega)\\
 &=c_0\sum_{\omega\in\Omega}p(\omega)
     +\sum_{i=1}^n c_i\sum_{\omega\in\Omega}X_i(\omega)p(\omega)\\
 &=c_0+\sum_{i=1}^n c_i\mathbb E[X_i].
\end{aligned}
\]
The same calculation applies when \(n=0\). No product of joint masses
was used. If \(X\leq Y\) pointwise, every term of
\(\sum_\omega(Y(\omega)-X(\omega))p(\omega)\) is nonnegative.
Linearity identifies that sum with \(\mathbb E[Y]-\mathbb E[X]\),
proving monotonicity.

For independent real-valued \(X,Y\), choose their finite ranges \(S,T\).
Apply the grouping formula to their joint variable and the function
\((x,y)\mapsto xy\). Independence and finite distributivity give
\[
\begin{aligned}
 \mathbb E[XY]
 &=\sum_{x\in S}\sum_{y\in T}xy\,p_{X,Y}(x,y)\\
 &=\sum_{x\in S}\sum_{y\in T}xy\,p_X(x)p_Y(y)\\
 &=\left(\sum_{x\in S}x p_X(x)\right)
   \left(\sum_{y\in T}y p_Y(y)\right)
 =\mathbb E[X]\mathbb E[Y].
\end{aligned}
\]
All sums are finite real sums, so no convergence or exchange-of-limit
assumption is involved.
\end{proof}
\end{annexitem}

\begin{annexitem}{annex:p15.finite-total-expectation}{thm:p15-finite-total-expectation}{Conditional\\means and their averaging property}
\begin{proof}
Partition the defining sum for \(\mathbb E[X]\) by the events \(H_i\).
For \(i\notin J\), every mass on \(H_i\) is zero, so this block
contributes zero. For \(i\in J\), multiply the definition of its
conditional expectation by \(\mathbb P(H_i)\). Thus
\[
 \mathbb E[X]
 =\sum_{i\in J}\sum_{\omega\in H_i}X(\omega)p(\omega)
 =\sum_{i\in J}\mathbb P(H_i)\mathbb E[X\mid H_i].
\]
The omitted terms are zero sums, not products of zero with undefined
conditional expectations.

Write \(Z=g(Y)\) as in the definition. For each \(y\in T\) with
\(p_Y(y)>0\), \(Z\) is constant at \(g(y)\) on \(\{Y=y\}\).
Consequently
\[
 \sum_{\omega:Y(\omega)=y}Z(\omega)p(\omega)
 =g(y)p_Y(y)
 =\sum_{\omega:Y(\omega)=y}X(\omega)p(\omega).
\]
For a null fiber, both sums are zero. Add these identities over
\(y\in D\) to obtain the stated averaging property; \(D=T\) gives
\(\mathbb E[Z]=\mathbb E[X]\).

Suppose \(W=h(Y)\) has the same averaging property. Take \(D=\{y\}\).
For \(p_Y(y)>0\), it follows that
\(h(y)p_Y(y)=g(y)p_Y(y)\), so \(h(y)=g(y)\). The remaining
fibers have probability zero and their union is null by finite additivity.
Hence \(\mathbb P(W=Z)=1\). This proves almost-sure uniqueness and also
shows why the chosen values on null fibers do not affect the result.
\end{proof}
\end{annexitem}
